A bank that uses a model to decide on credit, capital or liquidity must show, to its own validation function and to its supervisor, that the model does what it claims. The supervisory guide to internal models is the text under which such a validation function works \citep{ecb2025}. A validation test is a decision rule: it reads a sample, computes a statistic, and returns a light, a p-value or a flag. Like every decision rule it has two error rates. Its size is the probability that it raises an alarm when nothing is wrong, and its power is the probability that it raises one when a defect of a stated severity is present. A rule with a large size teaches the people who read it to ignore its alarms. A rule with a low power leaves the defect it was built for undetected. The literature that recommends a validation test states the statistic and, at most, a nominal level; it rarely states the size the rule actually has at the sample sizes in which it is used, and it almost never states what a correction of the rule costs in power.

This manuscript measures those two rates for three corrections of three tests in the validation battery of an open risk-model engine. The tests are the ones a machine-learning model of credit risk meets after deployment. The first is the coverage test of a split-conformal prediction set: the share of recent rows whose set holds the true label is compared with the target $1-\alpha$. The natural test, and the one the research behind the product prescribes, is a binomial test on the coverage indicators of an out-of-time window. The second is the resilience test: the loss of discrimination (the AUC) on the worst subsample of the evaluation rows, selected by how little they resemble the development rows, is compared with a fixed cut point on the AUC drop; the resilience test of the PiML toolbox evaluates performance degradation under out-of-distribution scenarios \citep{sudjianto2023}, and the fixed cut is the rule of the engine. The third is the situation-testing light of a fairness battery: the share of complainants of a protected group who are treated worse than comparable applicants of the other groups is compared with an amber cut point \citep{luong2011,alvarez2025}. Each of the three comparisons hides a source of error. The binomial coverage test ignores that the calibration quantile is itself random. The fixed AUC-drop cut ignores that the drop on the worst subsample is positive even when nothing has moved, because a subsample defined by distance from the development sample is a restricted region of the feature space. The situation-testing cut ignores that complainants share neighbours, so the count of significant complainants is far more variable than a binomial count.

\subsection*{The corrections and the measurement}

The proposed corrections are one for each test. The coverage test reads the covered count of the test window against the Beta-Binomial law that a calibration window of $n$ rows implies, and keeps the binomial p-value beside it. The resilience test reads the observed AUC drop against a permutation reference built from development rows only, with 199 draws by default. The situation-testing light sets its amber cut from a measured size and power instead of at 0.10: the engine's cut of 0.20 is the value that the rule of \cref{eq:cut-rule} selects on the engine's own design at 400 applicants, and the rule selects other cuts on the other three design and size pairs measured. None of the three corrections introduces a new statistic: each replaces a reference law that is correct only under an extra condition by the law that the rows actually have, and the measurement shows what that replacement changes in the rates. The proposed measurement is the part that is new to the validation practice: each rule and the rule it replaces are run on the same seeded draws of generators with a known truth, the rejection rate of each is reported with its Monte Carlo standard error, and, where the null law is known, the exact size is computed by summation, which has no sampling error.

\subsection*{Contributions}

The contributions are the following.
\begin{enumerate}
\item The size of the binomial coverage test is measured and derived. At a test window as long as the calibration window it rejects \CovSinBinB{} of exchangeable replications at the nominal level 0.05 (standard error \CovSeBinB{}), and at four times the calibration window \CovSinBinC{} (\CovSeBinC{}); the Beta-Binomial test rejects \CovSinBBB{} and \CovSinBBC{}, and its exact size never exceeds \CovBBMaxExact{} on a grid of \CovGridCells{} window pairs. A variance ratio $(m+n+1)/(n+2)$ explains the excess in closed form.
\item The power that the binomial test appears to have is shown to be its size. A binomial test whose cut is calibrated to the exact size 0.05 has the power of the Beta-Binomial test to within \ExPowGapMax{} in exact power, so the correction removes false alarms and almost no detections.
\item A permutation reference removes the false alarms of a fixed AUC-drop cut and pays for it in power, by an amount that shrinks as the evaluation sample grows; the size of a fixed cut depends on the evaluation sample size, so a cut chosen at one size does not carry to another.
\item The amber cut of the situation-testing light is shown to be a convention whose size depends on the number of applicants and on the data-generating process, and the cut 0.20 is measured to cost power against the cut 0.10 at small effects. Of the \HypTotal{} predictions, \HypRefuted{} are refuted, and the refutations are reported with the others.
\item The same functions the experiments call produce the committed evidence of the product: on the committed shift windows of a credit model the Beta-Binomial correction changes the light of one of five windows from red to amber, and the amber cut 0.20 sets \CmFairGreen{} of the \CmFairPairs{} mortgage year pairs of the committed fairness case green, where the cut 0.10 would set all \CmFairAboveTen{} pairs amber; the size of neither cut is measured at the committed number of complainants, and the green pairs have shares at least \CmFairGreenOverNull{} times the largest null mean share of the plants.
\end{enumerate}

\subsection*{Main result and structure}

The corrections remove most of the false alarms of the rules they replace. At 500 evaluation rows a fixed drop of 0.05 raises a false alarm in \ShRateFixedFiveCZ{} of the shift-free replications against \ShRatePermCZ{} for the permutation reference, which detects a planted shift of two standard deviations in \ShRatePermCT{} against \ShRateFixedFiveCT{} for the fixed cut. With 400 applicants the amber cut 0.20 raises a false alarm in \SiRateCutTwentySmallZ{} of the replications against \SiRateCutTenSmallZ{} for the cut 0.10, and detects a planted bias of gap 0.4 in \SiRateCutTwentySmallC{} against \SiRateCutTenSmallC{}. The cost in power is real and the experiments quantify it. The predictions that fail concern the amber cut of the situation light: its size depends on the number of applicants and on the data-generating process, and the manuscript reports those failures together with the confirmations.

\Cref{sec:related-work} places the work among the published methods, \cref{sec:problem} states the problem formally, \cref{sec:proposed-method} describes the three corrections and the measurement protocol, and \cref{sec:theory} derives the size of each rule where the null law is known. \Cref{sec:experimental-design} states the predictions, the data and the statistics; \cref{sec:results} reports the results, \cref{sec:ablation} the contribution of each component, and \cref{sec:discussion,sec:limitations} their meaning and their limits.
